%% ===================================================================
%% OverclusterDRG --- ALENEX submission (revised)
%% Compile: pdflatex -> bibtex -> pdflatex -> pdflatex
%% Requires siam.cls in the same directory.
%% ===================================================================
\documentclass{siam}

\usepackage{amsmath,amssymb}
\usepackage{algorithm}
\usepackage{algpseudocode}
\usepackage{booktabs}
\usepackage{graphicx}
\usepackage{subcaption}
\usepackage[hidelinks]{hyperref}

\graphicspath{
  {./}
  {./results/}
  {./results/multidata/}
  {./OverclusterDRG_Paper/}
  {./OverclusteringRKC/}
  {./overclustering_me/}
}

\newtheorem{theorem}{Theorem}
\newtheorem{lemma}{Lemma}
\newtheorem{corollary}{Corollary}
\newtheorem{definition}{Definition}
\newenvironment{conjecture}[1][]{\par\medskip\noindent
  \textbf{Conjecture (#1).}\itshape\ }{\par\medskip}
\makeatletter
\@ifundefined{proof}{%
  \newenvironment{proof}[1][Proof]{\noindent\textit{#1.}\ }%
    {\hfill$\square$\par\medskip}%
}{}
\makeatother
\newcommand{\remark}[1]{\par\medskip\noindent\textit{Remark.} #1\par\medskip}

\newcommand{\OPT}{\mathrm{OPT}}
\newcommand{\rz}{r_z}
\newcommand{\CL}{C_L}
\newcommand{\Copt}{C^{*}_{Q}}

\title{OverclusterDRG: A Deterministic, LP-Free Algorithm for
Robust $k$-Center with Outliers\thanks{Department of Computer
Science, University of Delhi. Supervised by Prof.\ Neelima Gupta.}}

\author{Rudra Bhardwaj}

\begin{document}
\maketitle

\begin{abstract}
The Robust $k$-Center problem asks for $k$ centers covering all but
$z$ outliers at minimum radius. Existing methods sit at two extremes:
the Charikar et al.\ greedy delivers quality at quadratic cost, while
the RKC-Gonzalez heuristic is fast but has no guarantee; the
randomised method of Ding et al.\ lies between them, but its trial
count grows linearly in $z$. We present OverclusterDRG, a
deterministic, LP-free, single-pass algorithm. The key idea is to run
Gonzalez for $k+O$ steps instead of $k$, building an overcomplete
pool, then select $k$ centers by greedy initialisation and $1$-swap
local search. On 36 configurations of UCI Adult ($N$ up to $20{,}000$),
OC-Local achieves a mean LP ratio of $1.147$, beating Charikar in
$33/36$ configurations (median $7\times$ faster) and Ding in $35/36$
(median $8\times$ faster). Every comparison is significant at
$p<10^{-5}$. The central theoretical result is an exact phase
transition: the pool provably contains a $3\cdot\OPT$ solution
whenever $O\ge z$, and contains no finite guarantee when $O<z$, so
$O=z$ is the smallest valid overcount. One extra Gonzalez step
additionally gives a free certified lower bound $\OPT\ge D_{k+z+1}/2$
at every run. The quality advantage is consistent across datasets:
$1.155$ on Diabetes Indicators and $1.171$ on Covertype against
interpolated LP bounds.
\end{abstract}

%% ===================================================================
\section{Introduction}
%% ===================================================================

A single distant point can dictate the entire $k$-center objective.
Robust $k$-Center (RKC) repairs this: given $N$ points, $k$ centers,
and an outlier budget $z$, find the smallest radius covering all but
$z$ points. The problem is NP-hard, and the outlier budget interacts
badly with farthest-point methods---the points to ignore are exactly
the ones that attract the algorithm.

Two baselines bracket the problem. Charikar et al.\ \cite{charikar2001}
give a $3$-approximation, but on UCI Adult it ran $5{,}528\times$ slower
than RKC-Gonzalez while improving the radius by only $8.32\%$. The
Gonzalez adaptation is fast and competitive in practice but has no
approximation guarantee---one outlier can hijack a center. Ding et
al.\ \cite{ding2019} interpolate with a randomised strategy requiring
$T=\lceil(z{+}1)\ln 100\rceil$ trials, so runtime grows with $z$ and
the output is random.

OverclusterDRG fills the gap: run Gonzalez for $k+O$ steps, building a
pool of $k+O$ candidates, then select the best $k$-subset. Outliers
still consume steps, but the $O$ extra steps let the pool absorb them
and still hold a representative for each true cluster. The right slack
is exactly $O=z$---this is a theorem, not a hyperparameter.

\noindent\textbf{Contributions.}
\begin{enumerate}
  \item \textbf{Algorithm.} OverclusterDRG, deterministic, LP-free,
  single-pass, running in $O(Nkz)$.
  \item \textbf{Theory.} An exact phase transition at $O=z$: the pool
  contains a $3\cdot\OPT$ solution for $O\ge z$, and the ratio is
  unbounded for $O<z$ (Theorems~\ref{thm:sufficiency} and
  \ref{thm:necessity}). One extra Gonzalez step gives a free certified
  lower bound $\OPT\ge D_{k+z+1}/2$ at every $N$.
  \item \textbf{Experiments.} Mean LP ratio $1.147$ on Adult (exact LP),
  beating Charikar in $33/36$ and Ding in $35/36$ configurations;
  $p<10^{-5}$ throughout. Results are consistent on Diabetes and
  Covertype.
\end{enumerate}

%% ===================================================================
\section{Problem Setup}
%% ===================================================================
\label{sec:setup}

\begin{definition}[Robust $k$-Center]
Given $(X,d)$ with $|X|=N$, integers $k\ge1$ and $z\ge0$, find
$C\subseteq X$ with $|C|=k$ and $S\subseteq X$ with $|S|\le z$
minimising
\[
  R(C,S) = \max_{x\in X\setminus S} d(x,C), \qquad
  \OPT = \min_{|C|=k,\,|S|\le z} R(C,S).
\]
\end{definition}

For a fixed $C$, the optimal $S$ is simply the $z$ farthest points,
so we write $\rz(C,X)$ for the $(z{+}1)$-th largest distance from
$X$ to $C$, and $\OPT=\min_{|C|=k}\rz(C,X)$.

\paragraph{LP lower bound.}
We binary-search candidate radii and solve a covering LP at each
probe; the smallest feasible radius $R_{\mathrm{LP}}$ satisfies
$R_{\mathrm{LP}}\le\OPT$. We report $\rz(C)/R_{\mathrm{LP}}$
throughout. The full LP formulation is in Appendix~\ref{app:lp}.

\paragraph{Two optima.}
Let $\OPT_{m,0}=\min_{|C|=m}\max_{x\in X}d(x,C)$ denote the
standard $m$-center optimum with no outliers. The relationship between
$\OPT_{k+O,0}$ and $\OPT_{k,z}$ is the hinge of the theory.

%% ===================================================================
\section{Related Work}
%% ===================================================================

Gonzalez \cite{gonzalez1985} gave the $2$-approximation for $k$-center
in $O(Nk)$; its RKC adaptation runs $\rz$ on the result but has no
guarantee.

Charikar et al.\ \cite{charikar2001} gave the first constant-factor
guarantee for outlier clustering; the RKC variant certifies $3r^{*}$
at quadratic cost. Ding et al.\ \cite{ding2019} achieve a good-quality
randomised method, but the trial count $T=O(z\log(1/\delta))$ limits
scalability.

Local search gives the best known bounds for $k$-median: Arya et al.\
\cite{arya2004} proved a $(3+2/p)$-approximation for $p$-swap, via
telescoping sums over swaps. For $k$-center the objective is a maximum,
not a sum, and the argument breaks entirely. Goyal and Jaiswal
\cite{goyal2021} give tight FPT approximations for constrained
$k$-center variants; these target a different regime.

No prior algorithm is simultaneously deterministic, LP-free, fast at
large $N$ and large $z$, and provably tied to $\OPT$.

%% ===================================================================
\section{The OverclusterDRG Algorithm}
%% ===================================================================
\label{sec:algorithm}

The algorithm overclusters in Phase~1 and then selects in Phase~2.

\paragraph{Phase 1: pool construction.}
Run Gonzalez for $k+O$ steps. Start from the point nearest to the
empirical centroid (making the run deterministic). At each step add
the point farthest from the current pool. The result is a pool
$Q=\{q_1,\dots,q_{k+O}\}$, built in $O(N(k{+}O))$ time. Default: $O=z$.

\paragraph{Phase 2a: forward greedy (OC-Forward).}
Pre-compute the $N\times(k{+}O)$ distance matrix. Greedily add pool
points, each time choosing the one that minimises $\rz$ of the current
set. Cost: $O(NkO)$.

\paragraph{Phase 2b: $1$-swap local search (OC-Local).}
Starting from the forward-greedy solution, scan all pairs
$(c_i, q)\in C\times(Q\setminus C)$ and accept the first swap that
strictly decreases $\rz$. Repeat until no swap improves. The number of
distinct subsets is finite and the objective strictly decreases each
step, so termination is guaranteed.

\begin{algorithm}[t]
\caption{OverclusterDRG$(X, k, z, O{=}z)$}
\label{alg:main}
\begin{algorithmic}[1]
\State $Q \gets \textsc{Gonzalez}(X,\, k+O)$
\State $C_0 \gets \textsc{ForwardGreedy}(X, Q, k, z)$
\State $C \gets \textsc{LocalSearch}(X, Q, C_0, z)$
\State \Return $C$
\end{algorithmic}
\end{algorithm}

\paragraph{Complexity.}
With $O=z$ the total work is $O(Nkz)$, dominated by Phase~2. Our
measured exponent is $O(N^{1.85})$ ($R^2=0.976$). Memory is
$O(Nz)$---about $7.9$\,GB at $N=99{,}000$, $k=20$, $z=10\%$, which can
be block-processed if needed.

%% ===================================================================
\section{Theoretical Analysis}
%% ===================================================================
\label{sec:theory}

The guarantee rests on two lemmas and two theorems.

\subsection{Why $O = z$ is the right overcount}

\begin{lemma}[Pool guarantee]
\label{lem:pool}
Let $Q=\textsc{Gonzalez}(X,k{+}O)$. There exists a $k$-subset
$C_0\subseteq Q$ with
$\rz(C_0,X)\le \OPT_{k,z} + 2\cdot\OPT_{k+O,0}$.
\end{lemma}

\begin{proof}
Let $(C^{*},S^{*})$ be an optimal RKC solution. For each
$c^{*}_j\in C^{*}$, pick the nearest pool point $q_j$; the Gonzalez
$2$-approximation gives $d(c^{*}_j,q_j)\le 2\cdot\OPT_{k+O,0}$. Set
$C_0=\{q_1,\dots,q_k\}$. For any inlier $x\notin S^{*}$ with nearest
optimal center $c^{*}_j$:
\[
d(x,C_0)\le d(x,c^{*}_j)+d(c^{*}_j,q_j)\le\OPT+2\cdot\OPT_{k+O,0}.
\]
Discarding $S^{*}$ gives $\rz(C_0)\le\OPT+2\cdot\OPT_{k+O,0}$.
\end{proof}

The bound mixes two different optima. The next lemma closes the gap
when $O\ge z$.

\begin{lemma}[OPT comparison]
\label{lem:optcompare}
For $O\ge z$:\; $\OPT_{k+O,0}\le\OPT_{k,z}$.
\end{lemma}

\begin{proof}
Take the $k$ optimal RKC centers and add one center at each of the
$z$ outliers. This covers all $N$ points within radius $\OPT$, so
$\OPT_{k+z,0}\le\OPT$. More centers can only help, so for $O\ge z$,
$\OPT_{k+O,0}\le\OPT_{k+z,0}\le\OPT$.
\end{proof}

\begin{theorem}[Sufficiency]
\label{thm:sufficiency}
For $O\ge z$, the pool contains a $k$-subset with $\rz\le 3\cdot\OPT$.
\end{theorem}

\begin{proof}
Lemma~\ref{lem:pool} gives $\rz(C_0)\le\OPT+2\cdot\OPT_{k+O,0}$.
Lemma~\ref{lem:optcompare} gives $\OPT_{k+O,0}\le\OPT$.
Substituting: $\rz(C_0)\le\OPT+2\cdot\OPT=3\cdot\OPT$.
\end{proof}

The converse shows $O=z$ is not just sufficient but the minimum.

\begin{theorem}[Necessity]
\label{thm:necessity}
For any $k\ge 2$, $z\ge 1$, $M>0$, there is an instance with
$\OPT_{k,z}=\varepsilon$ on which the pool with $O=z-1$ satisfies
$\min_{C\subseteq Q,|C|=k}\rz(C)\ge M\varepsilon$.
\end{theorem}

\begin{proof}[Proof sketch]
Place $k$ tight clusters of $m>z$ inliers each, and $z$ outliers at
distance $R\gg\varepsilon$ in orthogonal directions. The optimum covers
the $k$ clusters and discards the $z$ outliers at cost $\varepsilon$.
Since the outliers are the farthest points, Gonzalez selects all $z$
of them before any second cluster gets a representative. With $O=z-1$,
the pool has the starting inlier (from one cluster), the $z$ outliers,
and only $k-2$ further cluster representatives---leaving at least one
cluster with no pool point. Every $k$-subset of the pool either wastes
a slot on an outlier or leaves the orphaned cluster uncovered; since
that cluster has $m>z$ inliers, none of the $z$ budget slots can
discard them all. The resulting radius is at least $D-2\varepsilon\ge
M\varepsilon$ by choosing $D=2(M{+}1)\varepsilon$.
\end{proof}

\begin{corollary}[Exact phase transition]
\label{cor:phase}
$O=z$ is the smallest overcount with a finite pool guarantee. Below
$z$ the ratio is unbounded ($k\ge 2$); at and above $z$ the pool
always contains a $3\cdot\OPT$ subset.
\end{corollary}

The intuition: Gonzalez wastes exactly one step per outlier. With
$k+O$ total steps and $z$ wasted on outliers, the useful steps are
$k+O-z$. At $O=z$ this is exactly $k$---one representative per true
cluster. Below $z$, one cluster is left without a representative.

\subsection{A free certified lower bound}

\begin{lemma}[Certified lower bound]
\label{lem:cert}
If $N\ge k{+}z{+}1$, then $\OPT\ge D_{k+z+1}/2$, where $D_t$ is the
selection distance of the $t$-th Gonzalez point.
\end{lemma}

\begin{proof}
The $k+z+1$ Gonzalez points are pairwise at distance $\ge D_{k+z+1}$:
each $q_t$ was selected at distance $D_t\ge D_{k+z+1}$ from all
predecessors. The optimal solution discards at most $z$ of them, so
at least $k+1$ are inliers. By pigeonhole two share an optimal center,
giving $2\cdot\OPT\ge D_{k+z+1}$.
\end{proof}

Every run of OverclusterDRG automatically computes $D_{k+z+1}$ by
running one step past the pool. This yields the per-instance certificate
\[
  \frac{\rz(C)}{\OPT} \;\le\; \frac{2\,\rz(C)}{D_{k+z+1}}
\]
with no LP. On real data the certified ratio averages $3.6$ (Adult),
$3.1$ (Diabetes), $3.3$ (Covertype)---valid but conservative. The
bound is loosest when $z$ is large relative to $N$: running
$k+z+1\gg k$ Gonzalez steps drives $D_{k+z+1}$ well below $\OPT$.
The LP is about $3\times$ tighter where available. We recommend the
certificate as a quick sanity bound and the LP for ratio-accurate
comparison.

\subsection{Local search: what Phase~2 finds}

Exhaustive search over all $k$-subsets of $Q$ retrieves the guaranteed
$3\cdot\OPT$ subset (Theorem~\ref{thm:sufficiency}), but costs
$O(N\binom{k+O}{k})$---astronomical at $O=z$. The practical
$1$-swap local search is cheaper but is not exact:

\begin{theorem}[$1$-swap can fail]
\label{thm:neg}
For any $D>0$, there is an instance where a greedy-initialised $1$-swap
local optimum on $Q$ satisfies $\rz(\CL)/\OPT>D$.
\end{theorem}

\begin{proof}[Proof sketch]
Take $k$ equal-sized clusters and two singletons at slightly greater
distance. The optimum covers the clusters and discards the singletons.
Because the singletons are the two farthest points, the Gonzalez pool
places them first; forward greedy inserts both before enough cluster
representatives, leaving two clusters orphaned. No single swap covers
both at once, so the local optimum is stuck. The radius is the
inter-cluster distance $D'$; by choosing $D'$ large, the ratio
$D'/\varepsilon$ is unbounded.
\end{proof}

This failure needs adversarially equal inter-cluster distances, a
structure real datasets rarely produce. On $400$ small random
subsamples of Adult ($N\le 1{,}000$, $k\in\{5,10\}$), $1$-swap matched
exhaustive search exactly in $97.2\%$ of instances; the maximum
observed gap was $4.82\%$.

We conjecture that under generic (pairwise-distinct) distances and with
each optimal cluster large enough, $1$-swap satisfies
$\rz(\CL)\le(1+2/k)\cdot f(O)$ where $f(O)$ is the pool optimal.
This is supported by every one of the $400$ test instances but remains
open.

%% ===================================================================
\section{Experiments}
%% ===================================================================
\label{sec:experiments}

\subsection{Setup}
\label{sec:expsetup}

\begin{table}[t]
\caption{Datasets. All features z-score standardised.}
\label{tab:datasets}
\centering
\begin{tabular}{lrrl}
\toprule
Dataset & Rows used & Features & Source \\
\midrule
UCI Adult & 32{,}561 & 6 & \cite{adult_dataset} \\
Diabetes Indicators & 99{,}492 & 8 & UCI \\
Covertype & 30{,}000 of 581{,}012 & 10 & \cite{covertype_dataset} \\
HTRU2 & 17{,}898 & 8 & \cite{htru2} \\
\bottomrule
\end{tabular}
\end{table}

Table~\ref{tab:datasets} lists the datasets. The Adult grid uses
$N\in\{500,1000,2000,5000,10000,20000\}$, $k\in\{10,20,50\}$, and
outlier rate $z/N\in\{10\%,20\%\}$---36 configurations total. Diabetes
substitutes $N=15{,}000$ for $N=2{,}000$ (36 configurations); Covertype
uses $N\in\{500,1000,5000,10000,15000\}$ (30 configurations). Five
algorithms are compared: RKC-Gonzalez, OC-Forward, OC-Local, Charikar
et al.\ \cite{charikar2001} at its achieved radius, and Ding et al.\
\cite{ding2019} at $T=\lceil(z{+}1)\ln 100\rceil$ trials.

\paragraph{Preprocessing.}
Adult uses its six numeric features (age, fnlwgt, education-num,
capital-gain, capital-loss, hours-per-week) after z-scoring; categorical
columns are dropped. Covertype is subsampled to $30{,}000$ rows using
NumPy \texttt{default\_rng} with seed $42$, then z-scored. We use the
same subsample for all Covertype configurations.

\paragraph{LP bounds.}
All LP bounds are exact (Gurobi~12.0.3, feasibility tolerance
$10^{-6}$) on Adult---every configuration, at up to 94 hours for
$N=20{,}000$. On Diabetes and Covertype, the LP is too expensive at
large $N$, so we interpolate using the ratio learned from Adult
configurations; this is disclosed in all tables that use interpolated
values.

\paragraph{Certified lower bounds.}
As noted in Lemma~\ref{lem:cert}, every run produces a free lower
bound $\OPT\ge D_{k+z+1}/2$. The certified ratio for OC-Local averages
$3.6$ on Adult, $3.1$ on Diabetes, and $3.3$ on Covertype. The
certificate captures $36$--$39\%$ of the LP bound---it is a useful
sanity check but not a substitute for the LP.

\subsection{UCI Adult: full comparison}

\begin{table}[t]
\caption{LP ratios and runtimes across 36 UCI Adult configurations
(exact LP). Lower ratio is better; $1.000$ is optimal.}
\label{tab:adult}
\centering
\begin{tabular}{lcccccc}
\toprule
Algorithm & Mean & Std & Min & Max & Med.\ time (s) & Max time (s) \\
\midrule
LP oracle   & 1.000 & ---   & 1.000 & 1.000 & 4{,}586 & 338{,}755 \\
\textbf{OC-Local}  & \textbf{1.147} & 0.047 & 1.088 & 1.265 & 0.36 & 57 \\
OC-Forward  & 1.158 & 0.047 & 1.104 & 1.283 & 0.26 & 52 \\
Charikar    & 1.173 & 0.056 & 1.104 & 1.306 & 3.60 & 564 \\
Ding        & 1.235 & 0.032 & 1.179 & 1.331 & 2.22 & 431 \\
RKC-Gonzalez & 1.280 & 0.049 & 1.197 & 1.370 & 0.002 & 0.02 \\
\bottomrule
\end{tabular}
\end{table}

Table~\ref{tab:adult} gives the headline numbers. OC-Local achieves
the best mean ratio at $1.147$, ahead of Charikar ($1.173$), Ding
($1.235$), and RKC-Gonzalez ($1.280$). It is also faster: median
$0.36$\,s against Charikar's $3.60$\,s and Ding's $2.22$\,s.
OC-Local wins $33/36$ per-configuration comparisons against Charikar
at a median $7\times$ speedup, and $35/36$ against Ding at a median
$8\times$ speedup. The LP row is shown only for reference---its median
solve time is $4{,}586$ seconds.

\begin{figure}[t]
\centering
\includegraphics[width=0.95\linewidth]{fig3_quality_speed.png}
\caption{Quality--runtime tradeoff on UCI Adult. Each point is one
$(N,k,z/N)$ configuration; lower-left is better. OC-Local (green
diamonds) occupies the quality frontier at every runtime; no other
practical algorithm matches its radius at comparable speed.}
\label{fig:qualityspeed}
\end{figure}

\begin{figure}[t]
\centering
\includegraphics[width=0.95\linewidth]{fig1_ratio_vs_N.png}
\caption{LP ratio versus $N$, averaged over $k$, one panel per outlier
rate. OC-Local's ratio stays near $1.13$--$1.16$ throughout, while the
gap to Gonzalez and Ding grows with $N$.}
\label{fig:ratiovsN}
\end{figure}

Figure~\ref{fig:qualityspeed} shows the full quality--runtime picture:
OC-Local consistently dominates the lower-left corner. The advantage
is stable with dataset size (Figure~\ref{fig:ratiovsN}): OC-Local's
ratio stays between $1.13$ and $1.16$ from $N=500$ to $N=20{,}000$.

\begin{table}[t]
\caption{Paired tests across 36 Adult configurations.
Positive $\bar{d}$ means OC-Local is better.
Wilcoxon gives the signed-rank test.}
\label{tab:ttests}
\centering
\begin{tabular}{lcccc}
\toprule
Comparison & $\bar d$ & $t$ & $p$ & Wilcoxon $p$ \\
\midrule
OC-Local vs RKC-Gonzalez & $+0.134$ & $+13.3$ & $<10^{-14}$ & $2.9\times10^{-11}$ \\
OC-Local vs Ding         & $+0.089$ & $+14.3$ & $<10^{-15}$ & $5.8\times10^{-11}$ \\
OC-Local vs Charikar     & $+0.026$ & $+5.7$  & $<10^{-5}$  & $3.1\times10^{-8}$ \\
OC-Local vs OC-Forward   & $+0.012$ & $+7.9$  & $<10^{-8}$  & $2.9\times10^{-11}$ \\
\bottomrule
\end{tabular}
\end{table}

Table~\ref{tab:ttests} confirms none of the gaps are noise. Even the
smallest comparison---OC-Local vs its own forward-greedy
initialisation---is significant at $p<10^{-8}$: the $1$-swap phase
improves the mean ratio by $0.012$ in $36$ of $36$ configurations at a
median cost of $0.10$ seconds, so it should always be applied.
SimplifiedDRG (the poolless baseline) reaches only $1.295$; the pool
buys an $11.5\%$ improvement.

The $36$ configurations share a base dataset, so the paired $t$-tests
overstate independence. The Wilcoxon column (distribution-free) agrees
with every conclusion. As a further robustness check, restricting to
the $18$ smallest configurations ($N\le 2{,}000$, where the LP is
solved quickly) leaves every conclusion intact: OC-Local wins $17$--$18$
of $18$ against each baseline with Wilcoxon $p\le 1.3\times10^{-4}$.

\paragraph{Runtime scaling.}
Log-log regression gives OC-Local $\approx O(N^{1.85})$ ($R^2=0.976$)
and Charikar $\approx O(N^{2.12})$ ($R^2=0.999$). Ding's nominal
exponent is lower at $O(N^{1.10})$, but its trial count $(z{+}1)\ln100$
dominates: at $N=20{,}000$, $k=50$, $z=20\%$ Ding needs $18{,}423$
trials and takes $431$ seconds; OC-Local finishes in $57$ seconds.

\subsection{Local search validation}
\label{sec:localvalid}

Theorem~\ref{thm:neg} says $1$-swap can be arbitrarily bad. How often
does this actually happen? We ran OC-Local and exhaustive Phase~2
on $400$ small random subsamples of Adult ($N\le 1{,}000$,
$k\in\{5,10\}$, $z\le 10$)---the largest regime where exhaustive search
is feasible.

\begin{figure}[t]
\centering
\includegraphics[width=0.95\linewidth]{fig5_local_vs_exhaustive.png}
\caption{Gap between $1$-swap local search and exhaustive pool search
on $400$ instances. Local search is exact in $97.2\%$ of cases; the
maximum gap is $4.82\%$. The 11 non-zero gaps are shown in the
right panel.}
\label{fig:localexhaustive}
\end{figure}

OC-Local matched the exhaustive solution in $389$ of $400$ instances
($97.2\%$). The maximum observed gap was $4.82\%$, at pool size $20$
($k=10$, $z=10$), where exhaustive search evaluates $184{,}756$ subsets
in $10$ seconds and OC-Local finishes in $3$\,ms. The $11$ failures all
had the adversarial structure from Theorem~\ref{thm:neg}---small
clusters at near-uniform distances---which random real-data slices
essentially never produce. No observed gap exceeded $1+2/k$, which
supports the conjecture from Section~\ref{sec:theory}.

\subsection{Pool size analysis: the $O=z$ threshold}
\label{sec:oablation}

Corollary~\ref{cor:phase} says the transition happens exactly at $O=z$.
We test this by sweeping $O$ over $72$ configurations on Diabetes and
Covertype and measuring the pool-optimal radius as a function of $O$.

\begin{figure}[t]
\centering
\includegraphics[width=0.95\linewidth]{fig_O_ablation.png}
\caption{Radius (normalised to its value at $O=z$) versus overcount
$O$ on a log axis, across all $72$ configurations. Radius drops
steeply while $O<z$ and flattens near the threshold.}
\label{fig:oablation}
\end{figure}

\begin{figure}[t]
\centering
\begin{subfigure}{0.48\linewidth}
  \includegraphics[width=\linewidth]{fig_phase_transition.png}
  \caption{}
\end{subfigure}
\hfill
\begin{subfigure}{0.48\linewidth}
  \includegraphics[width=\linewidth]{fig_O_vs_zk.png}
  \caption{}
\end{subfigure}
\caption{(a) Radius vs $O$ for three representative configurations
with the $O=z$ threshold marked. (b) Empirical knee $O_{\rm knee}/z$
vs $z/k$ across 72 configurations (Pearson $r=-0.60$).}
\label{fig:phasepair}
\end{figure}

Figure~\ref{fig:oablation} shows the sweep: quality falls steeply left
of $O=z$ in every configuration, exactly as the theory predicts. Past
the threshold, quality keeps improving modestly. The empirical knee
$O_{\rm knee}$ (the smallest $O$ within $2\%$ of best pool quality)
sits between $1.0z$ and $1.7z$, with mean $1.55z$.

It is worth being precise about what the theory covers and what it
does not. Corollary~\ref{cor:phase} is about what the pool
\emph{contains}---the $3\cdot\OPT$ guarantee switches on at $O=z$.
The knee measures what OC-Local \emph{finds}, which keeps improving
past $z$. These are different questions, and the gap between them is
open.

\paragraph{Adaptive stopping.}
One can stop Gonzalez when $D_{k+t+1}\le\varepsilon D_{k+z+1}$,
obtaining a pool with worst-case factor $(1+2\varepsilon)\cdot\OPT$
(from the anytime bound). We ran this at
$\varepsilon\in\{1.0,0.5,0.25,0.1\}$. As expected, $\varepsilon=1.0$
stops exactly at $O=z$. Tighter $\varepsilon$ buys a better bound at
the cost of a larger pool: mean stopping at $3.7z$ ($\varepsilon=0.5$),
$6.0z$ ($0.25$), and $6.9z$ ($0.1$).

\begin{table}[t]
\caption{Pool-size settings across $72$ configurations. Gap is
relative to the best quality over all tested $O$.}
\label{tab:osettings}
\centering
\begin{tabular}{lcccc}
\toprule
Setting & Within 5\% & Mean gap & Max gap & Bound \\
\midrule
$O=z$        & 53/72 & $+3.7\%$ & $+11.8\%$ & $3\cdot\OPT$ \\
$O=2z$       & 69/72 & $+1.3\%$ & $+8.2\%$  & $<3\cdot\OPT$ \\
Adaptive ($\varepsilon=0.5$, $O\approx3.7z$) & \textbf{72/72} & ${<}1\%$ & ${<}3\%$ & $2\cdot\OPT$ \\
\bottomrule
\end{tabular}
\end{table}

Table~\ref{tab:osettings} summarises the tradeoff. $O=z$ is the right
default: 53/72 configurations land within $5\%$ of best. $O=2z$ raises
that to 69/72 at $2\times$ the runtime; the adaptive rule at
$\varepsilon=0.5$ reaches 72/72 with a tighter worst-case bound but
roughly $3.7\times$ the pool.

\paragraph{What predicts the knee.}
Pearson correlations with $O_{\rm knee}/z$ ($n=72$): $z/k$ has
$r=-0.60$ ($p<10^{-9}$), raw $z$ has $r=-0.61$, $N$ has $r=-0.59$,
and dimensionality is uncorrelated ($r=+0.04$, $p=0.72$). Under a
joint OLS fit ($R^2=0.43$), $z/k$ is the strongest standardised
predictor ($-0.25$). When each center covers many outliers, the knee
sits close to $z$; when outliers are few per center, slightly more
pool helps.

\begin{figure}[t]
\centering
\includegraphics[width=0.95\linewidth]{fig_synthetic.png}
\caption{Controlled synthetic sweeps of $z/k$, cluster separation, and
dimensionality. The knee stays near $z$ throughout; only $d\ge12$
pushes it above $2.5z$.}
\label{fig:synthetic}
\end{figure}

Synthetic experiments (Figure~\ref{fig:synthetic}) confirm this across
two decades of $z/k$, a range of cluster separations, and dimensions
$2$ to $20$.

\subsection{Cross-dataset validation}
\label{sec:multidata}

\begin{figure}[t]
\centering
\includegraphics[width=0.95\linewidth]{fig_ratio_both.png}
\caption{LP ratio versus $N$ on Diabetes (left) and Covertype (right).
The algorithm ordering and magnitudes replicate the Adult results.}
\label{fig:ratioboth}
\end{figure}

\begin{table}[t]
\caption{LP ratios on secondary datasets. Values are against
interpolated LP bounds. Charikar at achieved radius.}
\label{tab:multidata}
\centering
\begin{tabular}{llcccc}
\toprule
Dataset & Algorithm & Mean & Std & Min & Max \\
\midrule
Diabetes ($d{=}8$, $n{=}36$)
  & RKC-Gonzalez         & 1.923 & 0.256 & 1.449 & 2.414 \\
  & Ding                 & 1.227 & 0.048 & 1.143 & 1.330 \\
  & Charikar             & 1.201 & 0.044 & 1.145 & 1.313 \\
  & OC-Local             & \textbf{1.155} & 0.043 & 1.093 & 1.275 \\
Covertype ($d{=}10$, $n{=}30$)
  & RKC-Gonzalez         & 1.671 & 0.173 & 1.418 & 1.960 \\
  & Ding                 & 1.213 & 0.049 & 1.121 & 1.300 \\
  & Charikar             & 1.227 & 0.060 & 1.134 & 1.374 \\
  & OC-Local             & \textbf{1.171} & 0.051 & 1.095 & 1.294 \\
\bottomrule
\end{tabular}
\end{table}

Table~\ref{tab:multidata} shows the secondary results. OC-Local has
the best mean LP ratio on both datasets: $1.155$ on Diabetes and
$1.171$ on Covertype. Notably, the plain Gonzalez baseline degrades
sharply on these datasets ($1.92$ and $1.67$), which widens OC-Local's
margin relative to Adult. Figure~\ref{fig:ratioboth} shows the per-$N$
curves mirror the Adult pattern. On HTRU2, raw radii tell the same
story at $N=17{,}000$, $k=10$, $z=10\%$: OC-Local $1.79$ in $17$\,s
against Gonzalez $4.21$ and Charikar's achieved radius above OC-Local's.

\begin{table}[t]
\caption{OC-Local speedup over Ding at full certificate trial count.
OC-Local is faster in all 66 configurations.}
\label{tab:speedup}
\centering
\begin{tabular}{lcccc}
\toprule
Dataset & Mean & Median & Min & Max \\
\midrule
Diabetes  ($n{=}36$) & $2.0\times$ & $2.0\times$ & $1.3\times$ & $2.4\times$ \\
Covertype ($n{=}30$) & $2.6\times$ & $2.6\times$ & $2.1\times$ & $3.6\times$ \\
\bottomrule
\end{tabular}
\end{table}

OC-Local is faster than Ding in every one of the $66$ configurations,
by $2$--$3\times$ on average (Table~\ref{tab:speedup}). This is smaller
than the $7$--$8\times$ advantage on Adult because each Ding trial is
cheap and the trial count here is moderate. The more important
differences at scale are quality and determinism; speed is a consistent
but secondary advantage on these datasets.

\subsection{Scalability}
\label{sec:scalability}

\begin{table}[t]
\caption{OC-Local at large $N$ (Diabetes, $k=20$, $z=10\%$).
Entries marked $\sim$ are projected from the fitted scaling law.}
\label{tab:scale}
\centering
\begin{tabular}{rccc}
\toprule
$N$ & OC-Local radius & OC-Local time & Ding time \\
\midrule
20{,}000 & 2.060 & 57s & 133s \\
50{,}000 & $\sim$2.11 & $\sim$310s & $\sim$900s \\
99{,}000 & $\sim$2.15 & $\sim$1{,}100s & $\sim$3{,}800s \\
\bottomrule
\end{tabular}
\end{table}

Table~\ref{tab:scale} extends to the full Diabetes dataset. At
$N=99{,}000$, OC-Local is projected to finish in about $18$ minutes;
Ding's randomised certificate needs over an hour; Charikar's
$O(N^{2.12})$ scaling implies several hours; and the LP is out of
reach entirely. OC-Local is the only algorithm here that is both
deterministic and carries a formal statement about what it has
searched.

%% ===================================================================
\section{Conclusion}
%% ===================================================================

OverclusterDRG is a practical middle ground for Robust $k$-Center: fast
like the Gonzalez heuristic, quality like Charikar, and deterministic
throughout. The central contribution is the exact phase transition at
$O=z$: the pool contains a $3\cdot\OPT$ solution when the overcount
equals the outlier budget, and contains no finite guarantee below it.
The $1$-swap variant achieves mean LP ratios of $1.147$--$1.171$ across
three datasets, beats Charikar in $33/36$ and Ding in $35/36$
configurations, and remains the best practical option at $N\approx
10^{5}$ where the LP is unavailable.

Three problems remain open. First, when does $1$-swap always find the
good pool subset? We conjecture it does under generic distances and
large-enough clusters, but the proof requires a new argument for
max-type objectives---the standard $k$-median telescoping does not
apply. Second, is the $3\cdot\OPT$ pool bound tight? Our best
constructions give lower bounds near $2\cdot\OPT$. Third, does the
phase-transition analysis extend to $k$-median with outliers, where
the sum-type objective changes the geometry entirely?

\appendix
\section{The LP lower bound in full}
\label{app:lp}

For a candidate radius $r$, the feasibility LP has variables
$y_i\in[0,1]$ (fractional center mass at point $i$),
$s_j\in[0,1]$ (fractional outlier mass at $j$), and
$x_{ij}\ge 0$ defined only for pairs with $d(x_i,x_j)\le r$:
\[
\begin{aligned}
  \textstyle\sum_{i:\,d(x_i,x_j)\le r} x_{ij} + s_j &\ge 1
    &&\forall j\\
  x_{ij} &\le y_i &&\forall i,j\\
  \textstyle\sum_i y_i &\le k\\
  \textstyle\sum_j s_j &\le z
\end{aligned}
\]
The optimal integral solution at radius $\OPT$ is feasible for
$r=\OPT$, so $R_{\mathrm{LP}}\le\OPT$. We binary-search the $O(N^2)$
pairwise distances. Solver: Gurobi~12.0.3, primal/dual feasibility
tolerance $10^{-6}$. Adult LP times range from $100$--$1{,}096$\,s at
$N=2{,}000$ to $3.4\times10^{5}$\,s at $N=20{,}000$. On Diabetes and
Covertype the LP is interpolated from the solvable range; this is
disclosed in Section~\ref{sec:expsetup} and in Table~\ref{tab:multidata}.

\bibliographystyle{siam}
\bibliography{refs}

\end{document}
